% Part: first-order-logic
% Chapter: sequent-calculus
% Section: proof-theoretic-notions

\documentclass[../../../include/open-logic-section]{subfiles}

\begin{document}

\begin{editorial}
  In deze paragraaf staan de definities van bewijsbaarheid en
  consistentie voor natuurlijke deductie bij elkaar.
\end{editorial}

\iftag{FOL}
      {\olfileid{fol}{seq}{ptn}}
      {\olfileid{pl}{seq}{ptn}}

\olsection{Begrippen over formele afleidingen}

\begin{explain}
Eerder hebben we belangrijke semantische begrippen gedefinieerd:
geldigheid, semantisch gevolg en vervulbaarheid. Die begrippen gaan
over de vraag of !!{sentence}s waar zijn in !!{structure}s. Nu
definiëren we de bijbehorende \emph{bewijstheoretische begrippen}. Die
gaan over de !!{derivability} of !!{nonderivability} van bepaalde
sequenten. Een belangrijke ontdekking was dat de semantische en de
bewijstheoretische begrippen precies met elkaar overeenkomen. De
\emph{correctheidsstelling} en de \emph{volledigheidsstelling} zeggen
dat dit zo is.
\end{explain}

\begin{defn}[Stellingen]
!!^a{sentence}~$!A$ heet een \emph{stelling} als er in~$\Log{LK}$
!!a{derivation} van de sequent $\quad \Sequent !A$ bestaat. Is dat zo,
dan schrijven we $\Proves !A$. Is $!A$ geen stelling, dan schrijven we
$\Proves/ !A$.
\end{defn}

\begin{defn}[!!^{derivability}]
!!^a{sentence}~$!A$ is \emph{!!{derivable} uit} een verzameling
!!{sentence}s~$\Gamma$, genoteerd $\Gamma \Proves !A$, precies in het
volgende geval. Kies eerst een eindige deelverzameling~$\Gamma_0
\subseteq \Gamma$. Kies vervolgens een rij~$\Gamma_0'$ die uit de
!!{sentence}s in~$\Gamma_0$ bestaat. $\Log{LK}$ moet de sequent
$\Gamma_0' \Sequent !A$ !!{derive}s. Als $!A$ niet !!{derivable} is uit $\Gamma$, schrijven
we $\Gamma \Proves/ !A$.
\end{defn}

De regels voor contractie, verzwakking en verwisseling zorgen ervoor
dat de volgorde en het aantal !!{sentence}s in~$\Gamma_0'$ niet uitmaken.
Als $\Gamma_0' \Sequent !A$ !!{derivable} is, dan geldt dat ook voor
$\Gamma_0'' \Sequent !A$, zolang $\Gamma_0''$ dezelfde !!{sentence}s
bevat als~$\Gamma_0'$. Neem bijvoorbeeld $\Gamma_0 = \{!B, !C\}$. De
rijen $\Gamma_0' = \tuple{!B, !B, !C}$ en $\Gamma_0'' = \tuple{!C, !C,
!B}$ bevatten allebei alleen de !!{sentence}s uit~$\Gamma_0$. Als een
sequent met de ene rij !!{derivable} is, dan is de overeenkomstige
sequent met de andere rij dat ook. Dat blijkt bijvoorbeeld uit:
\begin{prooftree}
  \AxiomC{}
  \Deduce$!B, !B, !C \fCenter !A$
  \RightLabel{\LeftR{\Contraction}}
  \UnaryInf$!B, !C \fCenter !A$
  \RightLabel{\LeftR{\Exchange}}
  \UnaryInf$!C, !B \fCenter !A$
  \RightLabel{\LeftR{\Weakening}}
  \UnaryInf$!C, !C, !B \fCenter !A$
\end{prooftree}
Vanaf nu gebruiken we daarom een verkorte schrijfwijze. Als $\Gamma_0$
een eindige verzameling !!{sentence}s is, kan $\Gamma_0 \Sequent !A$
staan voor elke sequent waarvan de linkerkant een rij met
!!{sentence}s uit~$\Gamma_0$ is. Als het nodig is, rekenen we stappen
met contractie, verwisseling en verzwakking er stilzwijgend bij.

\begin{defn}[Consistentie]
Een verzameling zinnen~$\Gamma$ is \emph{inconsistent} als uit een
eindig deel~$\Gamma_0 \subseteq \Gamma$ een sequent met een lege
rechterkant kan worden afgeleid: $\Log{LK}$ !!{derive}s dan
$\Gamma_0 \Sequent \quad$. Is $\Gamma$ niet inconsistent, dan geldt
voor elk eindig deel $\Gamma_0 \subseteq \Gamma$ dat $\Log{LK}$ de
sequent $\Gamma_0 \Sequent \quad$ niet !!{derive}s. We noemen de
verzameling dan \emph{consistent}.
\end{defn}

\begin{prop}[Reflexiviteit]
\ollabel{prop:reflexivity}
Als $!A \in \Gamma$, dan $\Gamma \Proves !A$.
\end{prop}

\begin{proof}
De beginsequent $!A \Sequent !A$ is !!{derivable}. Bovendien geldt
$\{!A\} \subseteq \Gamma$.
\end{proof}

\begin{prop}[Monotonie]
\ollabel{prop:monotonicity}
Als $\Gamma \subseteq \Delta$ en $\Gamma \Proves !A$, dan $\Delta
\Proves !A$.
\end{prop}

\begin{proof}
Neem aan dat $\Gamma \Proves !A$. Dan is er een eindige $\Gamma_0
\subseteq \Gamma$ waarvoor $\Gamma_0 \Sequent !A$ !!{derivable} is.
Omdat $\Gamma \subseteq \Delta$, is diezelfde $\Gamma_0$ ook een
eindig deel van~$\Delta$. De !!{derivation} van $\Gamma_0 \Sequent !A$
laat dus ook zien dat $\Delta \Proves !A$.
\end{proof}

\begin{prop}[Transitiviteit]
\ollabel{prop:transitivity}
Als $\Gamma \Proves !A$ en $\{!A\} \cup \Delta \Proves !B$, dan
$\Gamma \cup \Delta \Proves !B$.
\end{prop}

\begin{proof}
Uit $\Gamma \Proves !A$ krijgen we een eindige $\Gamma_0 \subseteq
\Gamma$ en !!a{derivation}~$\pi_0$ van $\Gamma_0 \Sequent !A$. Uit
$\{!A\} \cup \Delta \Proves !B$ krijgen we een eindige $\Delta_0
\subseteq \Delta$ en !!a{derivation}~$\pi_1$ van $!A, \Delta_0
\Sequent !B$. Beschouw de volgende !!{derivation}, waarin we die twee
combineren:
\begin{prooftree}
  \AxiomC{}
  \RightLabel{$\pi_0$}
  \Deduce$\Gamma_0 \fCenter !A$
  \AxiomC{}
  \RightLabel{$\pi_1$}
  \Deduce$!A, \Delta_0 \fCenter !B$
  \RightLabel{\Cut}
  \BinaryInf$\Gamma_0, \Delta_0 \fCenter !B$
\end{prooftree}
Omdat $\Gamma_0 \cup \Delta_0 \subseteq \Gamma \cup \Delta$, laat
deze afleiding zien dat $\Gamma \cup \Delta \Proves !B$.
\end{proof}

Hieruit volgen twee bijzondere gevallen. Ten eerste: als $\Gamma
\Proves !A$ en $!A \Proves !B$, dan $\Gamma \Proves !B$. Ten tweede:
als $!A_1, \dots, !A_n \Proves !B$ en $\Gamma \Proves !A_i$ voor
elke~$i$, dan $\Gamma \Proves !B$.

\begin{prop}
\ollabel{prop:incons}
$\Gamma$ is inconsistent precies als $\Gamma \Proves {!A}$ voor elke
zin~$!A$.
\end{prop}

\begin{proof}
Opgave.
\end{proof}

\begin{prob}
Bewijs \olref[fol][seq][ptn]{prop:incons}.
\end{prob}

\begin{prop}[Compactheid]
  \ollabel{prop:proves-compact}
  \begin{enumerate}
  \item Als $\Gamma \Proves !A$, dan is er een eindig deel $\Gamma_0
    \subseteq \Gamma$ waarvoor $\Gamma_0 \Proves !A$.
  \item Als elk eindig deel van~$\Gamma$ consistent is, dan is ook de
    hele verzameling $\Gamma$ consistent.
  \end{enumerate}
\end{prop}

\begin{proof}
  \begin{enumerate}
    \item Uit $\Gamma \Proves !A$ volgt dat er een eindig deel
      $\Gamma_0 \subseteq \Gamma$ is waarvoor de sequent $\Gamma_0
      \Sequent !A$ !!a{derivation} heeft. Daarmee geldt meteen
      $\Gamma_0 \Proves !A$.
    \item Stel dat $\Gamma$ inconsistent is. Dan is er een eindig deel
      $\Gamma_0 \subseteq \Gamma$ waarvoor $\Log{LK}$ de sequent
      $\Gamma_0 \Sequent \quad$ !!{derive}s. Dat eindige deel
      $\Gamma_0$ van~$\Gamma$ is dus zelf inconsistent.
  \end{enumerate}
\end{proof}
\end{document}
